% Part: first-order-logic
% Chapter: sequent-calculus
% Section: proving-things-quant

\documentclass[../../../include/open-logic-section]{subfiles}

\begin{document}

\olfileid[hi]{fol}{seq}{prq}

\olsection{परिमाणकों वाली \usetoken{p}{derivation}}

\begin{ex}
एक $\Log{LK}$-!!{derivation} दीजिए जिसका अंतिम अनुक्रम $\lexists[x][\lnot !A(x)]
\Sequent \lnot \lforall[x][!A(x)]$ हो।

परिमाणकों के साथ काम करते समय यह सुनिश्चित करना होता है कि अभिलाक्षणिक-चर
शर्त का उल्लंघन न हो। इसके लिए कभी-कभी कुछ अनुमानों को लगाने का क्रम बदलकर
देखना पड़ता है। सामान्यतः अभिलाक्षणिक-चर शर्त के अधीन नियमों को पहले सँभालना
सहायक होता है (पूरे प्रमाण में वे नीचे आएँगे)। साथ ही आगे देखकर यह अनुमान
लगाना उपयोगी है कि प्रारंभिक अनुक्रम कैसा होगा। हमारे मामले में वह
$!A(a) \Sequent !A(a)$ जैसा होना चाहिए। इसका अर्थ है कि परिमाणक-नियमों को
``उल्टी दिशा'' में लगाते समय वही पद---जिसे हम $a$ कहेंगे---$\lforall$ और
$\lexists$, दोनों नियमों के लिए चुनना होगा। हर नियम के लिए अलग पद चुनने पर
$!A(a) \Sequent !A(b)$ जैसा अनुक्रम मिलता, जो निस्संदेह व्युत्पन्न नहीं है।

सामान्य रीति से आरंभ करते हुए लिखते हैं:
\begin{prooftree}
\AxiomC{}
\UnaryInf$\lexists[x][\lnot !A(x)] \fCenter \lnot \lforall[x][!A(x)]$
\end{prooftree}
हम \LeftR{\exists} नियम या \RightR{\lnot} नियम, दोनों में से कोई भी लगा सकते
हैं। चूँकि \LeftR{\exists} नियम अभिलाक्षणिक-चर शर्त के अधीन है, इसलिए उसे
बाद के बजाय अभी सँभालना उचित है; अतः पहले वही लगाते हैं।
\begin{prooftree}
\AxiomC{}
\UnaryInf$ \lnot !A(a) \fCenter \lnot \lforall[x][!A(x)]$
\RightLabel{\LeftR{\lexists}}
\UnaryInf$ \lexists[x][\lnot !A(x)] \fCenter \lnot \lforall[x][!A(x)]$
\end{prooftree}
\LeftR{\lnot} और \RightR{\lnot} नियमों को उल्टी दिशा में लगाने पर मिलता है:
\begin{prooftree}
\AxiomC{}
\UnaryInf$\lforall[x][!A(x)] \fCenter !A(a)$
\RightLabel{\LeftR{\lnot}}
\UnaryInf$\lnot !A(a), \lforall[x][!A(x)] \fCenter $
\RightLabel{\LeftR{\Exchange}}
\UnaryInf$\lforall[x][!A(x)], \lnot !A(a) \fCenter $
\RightLabel{\RightR{\lnot}}
\UnaryInf$ \lnot !A(a) \fCenter \lnot \lforall[x] !A(x)$
\RightLabel{\LeftR{\lexists}}
\UnaryInf$ \lexists[x] \lnot !A(x) \fCenter \lnot \lforall[x] !A(x)$
\end{prooftree}
अब हमारा एकमात्र विकल्प \LeftR{\forall} नियम लगाना है। यह नियम
अभिलाक्षणिक-चर प्रतिबंध के अधीन नहीं है, इसलिए यहाँ कोई बाधा नहीं है। याद
रखें, हम $!A(a) \Sequent !A(a)$ के रूप का प्रारंभिक अनुक्रम पाना चाहते हैं;
अतः नियम लगाते समय $a$ को $!A$ का तर्क चुनना चाहिए।
\begin{prooftree}
\Axiom$!A(a) \fCenter !A(a)$
\RightLabel{\LeftR{\lforall}}
\UnaryInf$\lforall[x][!A(x)] \fCenter !A(a)$
\RightLabel{\LeftR{\lnot}}
\UnaryInf$\lnot !A(a), \lforall[x][!A(x)] \fCenter $
\RightLabel{\LeftR{\Exchange}}
\UnaryInf$\lforall[x][!A(x)], \lnot !A(a) \fCenter $
\RightLabel{\RightR{\lnot}}
\UnaryInf$ \lnot !A(a) \fCenter \lnot \lforall[x][!A(x)]$
\RightLabel{\LeftR{\lexists}}
\UnaryInf$ \lexists[x][ \lnot !A(x)] \fCenter \lnot \lforall[x][!A(x)]$
\end{prooftree}
विशेषतः परिमाणकों के साथ काम करते समय इस बिंदु पर दोबारा जाँचना आवश्यक है
कि अभिलाक्षणिक-चर शर्त का उल्लंघन तो नहीं हुआ। हमने इस शर्त के अधीन केवल
\LeftR{\exists} नियम लगाया है, और अभिलाक्षणिक चर~$a$ उसके निचले अनुक्रम
(अंतिम अनुक्रम) में नहीं आता; अतः यह सही !!{derivation} है।
\end{ex}

\begin{prob}
निम्नलिखित अनुक्रमों की !!{derivation}s दीजिए:
\begin{enumerate}
\item $\Sequent (\lforall[x][!A(x)] \land \lforall[y][!B(y)]) \lif
\lforall[z][(!A(z) \land !B(z))]$.
\item $\Sequent (\lexists[x][!A(x)] \lor \lexists[y][!B(y)]) \lif
\lexists[z][(!A(z) \lor !B(z))]$.
\item $\lforall[x][(!A(x) \lif !B)] \Sequent \lexists[y][!A(y)] \lif !B$.
\item $\lforall[x][\lnot !A(x)] \Sequent \lnot\lexists[x][!A(x)]$.
\item $\Sequent \lnot\lexists[x][!A(x)] \lif \lforall[x][\lnot !A(x)]$.
\item $\Sequent \lnot\lexists[x][\lforall[y][((!A(x,y) \lif \lnot
!A(y,y)) \land (\lnot !A(y,y) \lif !A(x,y)))]]$.
\end{enumerate}
\end{prob}

\begin{prob}
निम्नलिखित अनुक्रमों की !!{derivation}s दीजिए:
\begin{enumerate}
\item $\Sequent \lnot\lforall[x][!A(x)] \lif \lexists[x][\lnot!A(x)]$.
\item $(\lforall[x][!A(x)] \lif !B) \Sequent \lexists[y][(!A(y) \lif !B)]$.
\item $\Sequent \lexists[x][(!A(x) \lif \lforall[y][!A(y)])]$.
\end{enumerate}
(इन सभी में $\RightR{\Contraction}$~नियम आवश्यक है।)
\end{prob}

\end{document}
